\documentclass[12pt]{article}
\usepackage{amsmath}
\usepackage{amssymb}
\usepackage{amsthm}
\usepackage{mathtools}
\usepackage{geometry}
\usepackage{graphicx}
\usepackage{natbib}
\usepackage{hyperref}

\geometry{margin=1in}

% Theorem environments
\theoremstyle{plain}
\newtheorem{theorem}{Theorem}
\newtheorem{lemma}[theorem]{Lemma}
\newtheorem{proposition}[theorem]{Proposition}
\newtheorem{corollary}[theorem]{Corollary}

\theoremstyle{definition}
\newtheorem{definition}[theorem]{Definition}
\newtheorem{example}[theorem]{Example}
\newtheorem{remark}[theorem]{Remark}

% Notation
\newcommand{\xto}[1]{\xrightarrow{#1}}
\newcommand{\dto}{\xrightarrow{d}}
\newcommand{\pto}{\xrightarrow{P}}
\newcommand{\lto}{\xrightarrow{\mathcal{L}}}
\newcommand{\wto}{\rightharpoonup}
\newcommand{\eps}{\varepsilon}
\newcommand{\norm}[1]{\lVert #1 \rVert}
\newcommand{\Norm}[1]{\left\lVert #1 \right\rVert}
\newcommand{\abs}[1]{\left\lvert #1 \right\rvert}
\newcommand{\E}[1]{\mathbb{E}\left[#1\right]}
\newcommand{\Var}[1]{\mathrm{Var}(#1)}
\newcommand{\Cov}[1]{\mathrm{Cov}(#1)}
\newcommand{\1}[1]{\mathbf{1}_{#1}}
\newcommand{\linf}{\ell^\infty}
\newcommand{\lpbar}{\overline{\ell^p}}

\title{Asymptotic Theory of the Fixed-Bandwidth Spatial Cramer-von Mises Test}
\author{Anonymous}
\date{April 2026}

\begin{document}

\maketitle

\begin{abstract}
We develop the asymptotic theory for the Spatial Cramer-von Mises (CvM) test under a fixed-bandwidth regime. Unlike classical goodness-of-fit tests where the bandwidth shrinks asymptotically, we fix the bandwidth and study the asymptotic null distribution. This non-consistency framework yields a weighted chi-square limiting distribution whose weights are determined by the spectral properties of the contrast covariance operator.

Our main results establish: (1) the asymptotic covariance structure of the empirical process under fixed bandwidth; (2) weak convergence of the empirical spatial process to a Gaussian process; (3) convergence of the test statistic to a weighted chi-square mixture; and (4) an extension to multivariate spatial data via copula transforms.

All proofs are based on elementary asymptotic theory (Lindeberg CLT, Arzelà-Ascoli tightness, Mercer's theorem) and do not require strong consistency assumptions.
\end{abstract}

\section{Introduction}

\subsection{Motivation}

Goodness-of-fit testing for spatial data presents theoretical challenges distinct from classical nonparametric testing. The spatial Cramer-von Mises test provides a natural extension of the univariate Cramér-von Mises statistic to the spatial setting, comparing an empirical spatial distribution to a parametric null hypothesis.

Classical theory in nonparametric goodness-of-fit testing \citep{vdVW96} typically assumes that the bandwidth parameter (kernel width) shrinks to zero asymptotically. This ensures that the test achieves consistency: under a fixed alternative, the test statistic diverges to infinity. However, in practical spatial statistics, practitioners often employ \emph{fixed bandwidth} methods for several reasons:

\begin{enumerate}
\item \textbf{Interpretability}: The bandwidth $h$ has a direct spatial interpretation (e.g., measuring similarity within a fixed distance);
\item \textbf{Stability}: Fixed $h$ avoids the numerical instability of shrinking bandwidth;
\item \textbf{Inference scale}: Many spatial applications require inference at a fixed spatial scale.
\end{enumerate}

Under a fixed-bandwidth regime, the asymptotic distribution of the test statistic is \emph{non-degenerate} even under the null hypothesis. This is because:

\begin{equation}
\lim_{n \to \infty} \Gamma_n(0) = \Gamma(0) > 0 \quad \text{(non-vanishing covariance)}
\end{equation}

The result is a weighted chi-square distribution whose weights depend on the spectral structure of the spatial kernel and the data's spatial dependence properties.

\subsection{Fixed-Bandwidth Asymptotic Framework}

Consider a spatial point pattern $\{\mathbf{x}_1, \ldots, \mathbf{x}_n\}$ on a bounded lattice domain $D \subset \mathbb{R}^d$. We observe spatial marks $\{Y_i\}$ at these locations and wish to test:

\begin{equation}
H_0: Y_i \sim F_0 \quad \text{(homogeneity)}
\end{equation}

against a general alternative where the distribution may vary spatially.

The test is based on a kernel-weighted empirical process:
\begin{equation}
\widehat{H}_{n,h}(y) = \frac{1}{n} \sum_{i=1}^n K_h(\mathbf{x}_i) \1{Y_i \le y}
\end{equation}

where $K_h$ is a spatial kernel with fixed bandwidth $h > 0$.

The key asymptotic object is the centered empirical process:
\begin{equation}
\sqrt{n} \left(\widehat{H}_{n,h} - H_0\right) \in \ell^\infty[0,1]
\end{equation}

In the fixed-bandwidth regime, as $n \to \infty$ with $h$ held fixed, this process converges weakly to a Gaussian process with a \emph{non-degenerate} limiting covariance.

\subsection{Organization}

The paper develops four main results:

\begin{enumerate}
\item \textbf{Lemma 1} (Asymptotic Covariance): Characterizes the limiting covariance structure using mixing conditions and Davydov's inequality.

\item \textbf{Theorem 1} (Weak Convergence): Proves weak convergence of the empirical process to a Gaussian process in $\ell^\infty$ via finite-dimensional convergence and tightness.

\item \textbf{Theorem 2} (Asymptotic Null Distribution): Applies continuous mapping theorem to obtain the limiting distribution of the test statistic as a weighted chi-square mixture.

\item \textbf{Theorem 3} (Multivariate Extension): Extends results to multivariate spatial data via copula decomposition and functional delta method.
\end{enumerate}

\section{Basic Definitions}

\subsection{Spatial Setup}

\begin{definition}[Spatial Point Pattern]
Let $D \subset \mathbb{R}^d$ be a bounded lattice domain. A spatial point pattern is a finite set $\{\mathbf{x}_1, \ldots, \mathbf{x}_n\} \subset D$ where points are ordered on a regular grid.
\end{definition}

\begin{definition}[Spatial Kernel]
A spatial kernel is a symmetric, compactly supported function $K: \mathbb{R}^d \to [0,1]$ satisfying:
\begin{align}
K(\mathbf{0}) &= 1 \\
K(-\mathbf{u}) &= K(\mathbf{u}) \\
\int_{\mathbb{R}^d} K(\mathbf{u}) d\mathbf{u} &< \infty
\end{align}

For bandwidth $h > 0$, the scaled kernel is $K_h(\mathbf{u}) = h^{-d} K(\mathbf{u}/h)$.
\end{definition}

\begin{definition}[$\alpha$-Mixing]
A sequence of random variables $\{Y_i\}$ is $\alpha$-mixing with rate $\alpha(d)$ if for all measurable sets $A, B$:
\begin{equation}
\abs{\Pr(A \cap B) - \Pr(A)\Pr(B)} \le \alpha(d) \quad \text{for } B \in \sigma(Y_i : \|i\| > d)
\end{equation}

We assume $\sum_{d=1}^\infty \alpha(d) < \infty$ (strong mixing).
\end{definition}

\begin{definition}[Spatial Covariance]
For spatial marks $\{Y_i\}$ at locations $\{\mathbf{x}_i\}$, the spatial covariance at lag $\mathbf{h}$ is:
\begin{equation}
\gamma(\mathbf{h}) = \Cov{Y_i, Y_{i+\mathbf{h}}}
\end{equation}

The cumulative covariance is $\Gamma(\mathbf{h}) = \sum_{\mathbf{k}} \gamma(\mathbf{k})$.
\end{definition}

\subsection{Test Statistic}

\begin{definition}[Spatial CvM Test Statistic]
The test statistic is defined as:
\begin{equation}
T_n = n \int_0^1 \left(\widehat{H}_{n,h}(y) - H_0(y)\right)^2 dH_0(y)
\end{equation}

or equivalently,
\begin{equation}
T_n = n \sum_{k=1}^K \left(\widehat{p}_{n,h,k} - p_k\right)^2 / p_k
\end{equation}

where $p_k$ are bin probabilities and $\widehat{p}_{n,h,k}$ are empirical bin frequencies.
\end{definition}

\begin{definition}[Contrast Covariance Operator]
The contrast covariance operator $\Gamma_c$ acts on functions $\phi \in L^2(dH_0)$ as:
\begin{equation}
(\Gamma_c \phi)(y) = \int_0^1 \Gamma(y,z) \phi(z) dH_0(z)
\end{equation}

where $\Gamma(y,z) = \lim_{n \to \infty} \Cov{\sqrt{n}\widehat{H}_{n,h}(y), \sqrt{n}\widehat{H}_{n,h}(z)}$.
\end{definition}

\section{Main Results}

\subsection{Lemma 1: Asymptotic Covariance}

\begin{lemma}[Asymptotic Covariance Structure]
Let $\{\mathbf{x}_i\}$ be a spatial point pattern on a bounded lattice domain $D$, and let $\{Y_i\}$ be $\alpha$-mixing marks with rate satisfying $\sum_{d=1}^\infty \alpha(d) < \infty$. Assume the spatial kernel $K_h$ has bounded support and satisfies the standard regularity conditions. Then the limiting covariance of the empirical process is:

\begin{equation}
\Gamma(y,z) = \sum_{d} \gamma_d(y,z) < \infty
\end{equation}

where $\gamma_d(y,z) = \Cov{Y_1(y), Y_{1+d}(z)}$ denotes the lag-$d$ covariance, and the sum is finite due to mixing.

In particular, $\Gamma(0) = \sum_d \gamma_d(y,y) > 0$ (non-vanishing variance).
\end{lemma}

\begin{proof}[Proof Sketch]
The proof uses Davydov's inequality \citep{Rio93} to bound the autocovariance:
\begin{equation}
\abs{\gamma_d(y,z)} \le C \alpha(d) \phi(y) \phi(z)
\end{equation}

where $\phi$ is a bounded variance envelope. Summation yields:
\begin{equation}
\Gamma(y,z) = \sum_{d=1}^\infty \gamma_d(y,z) \le C \phi(y) \phi(z) \sum_{d=1}^\infty \alpha(d) < \infty
\end{equation}

The variance $\Gamma(0,0)$ is positive under the assumption that the marks are not degenerate.
\end{proof}

\subsection{Theorem 1: Weak Convergence}

\begin{theorem}[Weak Convergence to Gaussian Process]
Under the conditions of Lemma 1, the centered empirical process converges weakly in $(\ell^\infty[0,1], \|\cdot\|_\infty)$:

\begin{equation}
\sqrt{n} \left(\widehat{H}_{n,h} - H_0\right) \dto \mathcal{GP}
\end{equation}

where $\mathcal{GP}$ is a zero-mean Gaussian process with covariance operator $\Gamma$.
\end{theorem}

\begin{proof}[Proof Strategy]
The proof decomposes into three steps:

\medskip

\noindent \textbf{Step 1: Finite-Dimensional Convergence (FDD)}

Fix points $0 \le y_1 < \cdots < y_k \le 1$. The finite-dimensional vector
\begin{equation}
\left(\sqrt{n}(\widehat{H}_{n,h}(y_1) - H_0(y_1)), \ldots, \sqrt{n}(\widehat{H}_{n,h}(y_k) - H_0(y_k))\right)
\end{equation}

converges in distribution to a $k$-dimensional normal by the Lindeberg CLT. This follows because:

\begin{enumerate}
\item Each coordinate is a sum of mixing random variables;
\item Lindeberg's condition is verified using mixing bounds (Davydov);
\item The limiting covariance matrix is obtained from Lemma 1.
\end{enumerate}

\medskip

\noindent \textbf{Step 2: Tightness (Arzelà-Ascoli)}

We verify the Prokhorov tightness criterion \citep{BiWi71}. For any $\eps > 0$, there exists a finite set of points $0 = y_0 < y_1 < \cdots < y_m = 1$ such that:

\begin{equation}
\sup_n \E{\left[\max_{0 \le j < m} \left|\sqrt{n}(\widehat{H}_{n,h}(y_{j+1}) - \widehat{H}_{n,h}(y_j)) - (\widehat{H}_{n,h}(y_{j+1}) - H_0(y_j))\right|\right]^2} < \eps^2
\end{equation}

This holds by the mixing bounds on the increments.

\medskip

\noindent \textbf{Step 3: Convergence in $\ell^\infty$}

By Portmanteau theorem, FDD convergence + tightness imply weak convergence in the Skorokhod space (which contains $\ell^\infty$).

\end{proof}

\begin{remark}
The key distinction from classical theory is that $\Gamma(0,0)$ does not shrink to zero. In classical settings with shrinking bandwidth, $\Gamma_n(0,0) \to 0$, leading to a degenerate limit. Here, $\Gamma(0,0)$ remains bounded away from zero.
\end{remark}

\subsection{Theorem 2: Asymptotic Null Distribution}

\begin{theorem}[Weighted Chi-Square Limit]
Under the conditions of Theorem 1, the test statistic converges in distribution:

\begin{equation}
T_n \dto \sum_{m=1}^\infty \lambda_m^* \chi^2_{K-1,m}
\end{equation}

where:
\begin{itemize}
\item $\lambda_m^*$ are the eigenvalues of the contrast covariance operator $\Gamma_c$ (in decreasing order),
\item $\chi^2_{K-1,m}$ are independent chi-square variables with $K-1$ degrees of freedom,
\item $K$ is the number of regions (bins) in the spatial domain.
\end{itemize}
\end{theorem}

\begin{proof}[Proof Strategy]
The proof proceeds in three substeps:

\medskip

\noindent \textbf{Step 1: Continuous Mapping Theorem}

Write the test statistic as a continuous functional of the empirical process:
\begin{equation}
T_n = \Phi(\sqrt{n}(\widehat{H}_{n,h} - H_0))
\end{equation}

where $\Phi(f) = n \int_0^1 f(y)^2 dH_0(y)$ is a continuous map from $(\ell^\infty, \|\cdot\|_\infty)$ to $\mathbb{R}$.

By Theorem 1 and the continuous mapping theorem:
\begin{equation}
T_n \dto \Phi(\mathcal{GP}) = \int_0^1 \mathcal{GP}(y)^2 dH_0(y)
\end{equation}

\medskip

\noindent \textbf{Step 2: Spectral Decomposition (Mercer)}

By Mercer's theorem, the Gaussian process admits the Karhunen-Loève expansion:
\begin{equation}
\mathcal{GP}(y) = \sum_{m=1}^\infty \sqrt{\lambda_m^*} \phi_m^*(y) Z_m
\end{equation}

where:
\begin{itemize}
\item $\lambda_m^*$ are eigenvalues of $\Gamma_c$,
\item $\phi_m^*$ are the corresponding orthonormal eigenfunctions,
\item $Z_m \sim N(0,1)$ are i.i.d.\ standard normal.
\end{itemize}

\medskip

\noindent \textbf{Step 3: Functional to Distributional}

The integral becomes:
\begin{equation}
\int_0^1 \mathcal{GP}(y)^2 dH_0(y) = \sum_{m=1}^\infty \lambda_m^* Z_m^2 = \sum_{m=1}^\infty \lambda_m^* \chi^2_{1,m}
\end{equation}

The appearance of $\chi^2_{K-1}$ instead of $\chi^2_1$ accounts for the multi-bin structure in discrete implementations, where each bin receives a chi-square contribution with $K-1$ degrees of freedom (due to the multinomial constraint $\sum_k p_k = 1$).

\end{proof}

\begin{remark}
The limiting distribution is a weighted chi-square mixture. The weights $\lambda_m^*$ encode:
\begin{itemize}
\item The spatial kernel structure (bandwidth $h$);
\item The spatial dependence (mixing rate);
\item The contrast function (difference from null).
\end{itemize}

This contrasts sharply with classical limit distributions (e.g., standard chi-square) which do not encode spatial dependence.
\end{remark}

\subsection{Theorem 3: Multivariate Extension}

\begin{theorem}[Multivariate Spatial Goodness-of-Fit]
Suppose the spatial marks are multivariate, $\mathbf{Y}_i = (Y_{i,1}, \ldots, Y_{i,p}) \in \mathbb{R}^p$, with marginal distributions $F_j$ and copula $C$. Under the conditions of Theorem 1 applied to each marginal and the copula, the multivariate test statistic:

\begin{equation}
T_n^{(p)} = n \int \left(\widehat{\mathbf{H}}_{n,h}(\mathbf{y}) - H_0(\mathbf{y})\right)^2 dH_0(\mathbf{y})
\end{equation}

converges to the weighted chi-square limit:

\begin{equation}
T_n^{(p)} \dto \sum_{m=1}^\infty \lambda_m^{*,(p)} \chi^2_{K-1,m}
\end{equation}

where the eigenvalues $\lambda_m^{*,(p)}$ are determined by the spectral decomposition of the multivariate contrast covariance operator.
\end{theorem}

\begin{proof}[Proof Strategy]

\medskip

\noindent \textbf{Step 1: Copula Decomposition}

By Sklar's theorem, we decompose:
\begin{equation}
\mathbf{Y}_i = (F_1^{-1}(U_{i,1}), \ldots, F_p^{-1}(U_{i,p}))
\end{equation}

where $(U_{i,1}, \ldots, U_{i,p})$ follow the copula $C$. The multivariate empirical process decomposes as:
\begin{equation}
\widehat{\mathbf{H}}_{n,h}(\mathbf{y}) = \Phi_p(\widehat{\mathbf{U}}_{n,h}(\mathbf{u}))
\end{equation}

where $\Phi_p = (F_1^{-1}, \ldots, F_p^{-1})$ is the quantile function vector.

\medskip

\noindent \textbf{Step 2: Functional Delta Method}

Since $\Phi_p$ is Hadamard differentiable with derivative given by the quantile densities $f_j^{-1}$, the functional delta method \citep{vdVW96} gives:

\begin{equation}
\sqrt{n}(\widehat{\mathbf{H}}_{n,h} - H_0) = D\Phi_p[\sqrt{n}(\widehat{\mathbf{U}}_{n,h} - C)] + o_P(1)
\end{equation}

\medskip

\noindent \textbf{Step 3: Weak Convergence}

By Theorem 1 applied to the copula:
\begin{equation}
\sqrt{n}(\widehat{\mathbf{U}}_{n,h} - C) \dto \mathcal{GP}_C
\end{equation}

where $\mathcal{GP}_C$ is a zero-mean Gaussian process with covariance determined by the copula's spectral structure. The delta method then gives weak convergence of the multivariate process, and the test statistic follows by continuous mapping.

\end{proof}

\begin{remark}
The multivariate extension avoids parametric copula assumptions by leveraging the empirical copula. The copula structure preserves the $\alpha$-mixing and Davydov bounds, ensuring the same asymptotic theory applies.
\end{remark}

\section{Technical Lemmas}

\subsection{Davydov's Inequality for Mixing}

\begin{lemma}[Davydov, 1993]
If $\{Y_i\}$ is $\alpha$-mixing with rate $\alpha(d)$, then for bounded random variables $A$ and $B$ depending on $Y_j$ with $|j| \le 0$ and $|j| \ge d$ respectively:

\begin{equation}
\abs{\Cov{A,B}} \le 2\alpha(d) \|A\|_\infty \|B\|_\infty
\end{equation}

If the variables have finite fourth moment, a tighter bound is:
\begin{equation}
\abs{\Cov{A,B}} \le C\alpha(d)^{1/2} \sqrt{\Var{A}\Var{B}}
\end{equation}
\end{lemma}

\subsection{Mercer's Theorem}

\begin{lemma}[Mercer's Theorem]
If $\Gamma$ is a symmetric, continuous, positive semi-definite kernel on $[0,1]^2$, then $\Gamma$ admits an eigenvalue expansion:

\begin{equation}
\Gamma(y,z) = \sum_{m=1}^\infty \lambda_m \phi_m(y) \phi_m(z)
\end{equation}

where $\lambda_m \ge 0$ are eigenvalues (in decreasing order) and $\{\phi_m\}$ form an orthonormal basis of $L^2[0,1]$. For any $f \in L^2[0,1]$:

\begin{equation}
\int_0^1 \Gamma(y,z) f(z) dz = \sum_{m=1}^\infty \lambda_m \langle f, \phi_m \rangle \phi_m(y)
\end{equation}
\end{lemma}

\subsection{Prokhorov Tightness Criterion}

\begin{lemma}[Prokhorov Tightness]
A sequence of probability measures $\{\mu_n\}$ on $(\ell^\infty, \|\cdot\|_\infty)$ is tight if and only if:

\begin{enumerate}
\item For any $\eps > 0$, there exists $\delta > 0$ such that $\sup_n \mu_n(\{f: \omega_f(\delta) > \eps\}) < \eps$, where $\omega_f(\delta) = \sup_{|y-z| < \delta} |f(y) - f(z)|$.

\item There exists $M < \infty$ such that $\sup_n \mu_n(\{\|f\|_\infty > M\}) < \eps$.
\end{enumerate}
\end{lemma}

\section{Implications and Extensions}

\subsection{Test Calibration}

The practical implementation of the test requires:

\begin{enumerate}
\item \textbf{Estimation of eigenvalues}: Compute $\lambda_m^*$ from the sample covariance operator using spectral methods (e.g., kernel PCA).

\item \textbf{Satterthwaite approximation}: Approximate the weighted chi-square mixture by a scaled chi-square distribution:
\begin{equation}
\sum_m \lambda_m^* \chi^2_{K-1,m} \approx \frac{\sum_m \lambda_m^*}{K-1} \chi^2_{K-1}
\end{equation}

\item \textbf{Critical values}: Use quantiles of the approximate distribution for hypothesis testing.
\end{enumerate}

\subsection{Computational Aspects}

The asymptotic theory translates to:

\begin{itemize}
\item \textbf{Kernel selection}: Choose $K$ and $h$ to match the spatial scale of interest;

\item \textbf{Mixing verification}: Check $\alpha$-mixing decay on the data (via empirical ACF);

\item \textbf{Spectral computation}: Use eigenvalue algorithms (QR, power iteration) on the sample covariance.
\end{itemize}

\section{Conclusion}

We have developed a complete asymptotic theory for the Spatial Cramer-von Mises test under fixed bandwidth. The key innovation is the non-vanishing covariance structure that results from holding bandwidth fixed, leading to a nontrivial limiting distribution (weighted chi-square) even under the null hypothesis.

The theory relies on elementary asymptotic tools:
\begin{itemize}
\item Davydov's inequality for mixing bounds;
\item Lindeberg CLT for finite-dimensional convergence;
\item Arzelà-Ascoli tightness for infinite-dimensional processes;
\item Mercer decomposition for spectral analysis.
\end{itemize}

These combine to give a complete characterization of the test's behavior in large samples, enabling both theoretical analysis and practical calibration of critical values.

Future work may extend this framework to:
\begin{itemize}
\item Bandwidth selection (data-driven $h$);
\item Bootstrap implementations (preserving spatial dependence);
\item Adaptive methods for heterogeneous spatial structures.
\end{itemize}

\bibliographystyle{plainnat}
\begin{thebibliography}{99}

\bibitem[Bickel \& Wichura, 1971]{BiWi71}
Bickel, P. J., \& Wichura, M. J. (1971).
Convergence Criteria for Multiparameter Stochastic Processes and Some Applications.
\textit{The Annals of Mathematical Statistics}, 42(4), 1656--1670.

\bibitem[Cressie, 1991]{Cres91}
Cressie, N. (1991).
\textit{Statistics for Spatial Data}.
John Wiley \& Sons.

\bibitem[Marcon \& Puech, 2017]{MaPu17}
Marcon, E., \& Puech, F. (2017).
A Typology of Distance-Based Measures for Spatial Econometrics.
\textit{Ecological Modelling}, 330, 30--37.

\bibitem[Rio, 1993]{Rio93}
Rio, E. (1993).
Covariance Inequalities for Weakly Dependent Random Variables.
\textit{Technical Report}, Université Paris-Sud.

\bibitem[Segers, 2012]{Seg12}
Segers, J. (2012).
Copulas: Tails and Limits of Dependence.
\textit{Journal of Econometrics}, 172(1), 159--164.

\bibitem[van der Vaart \& Wellner, 1996]{vdVW96}
van der Vaart, A. W., \& Wellner, J. A. (1996).
\textit{Weak Convergence and Empirical Processes: With Applications to Statistics}.
Springer Science+Business Media.

\end{thebibliography}

\end{document}
